\section{Related Work and Theoretical Framing}

\subsection{From interface genealogy to pressability}

The button has a history. Plotnick's work shows that the push button emerged not just as a convenience feature but as a political and sensory technology that promised instantaneity while reorganizing who or what could act at a distance \citep{Plotnick2018,PlotnickInterface2012}. In parallel, media framings of ``push-button warfare'' turned remote activation into a moral drama of effortlessness, abstraction, and command \citep{PlotnickWarfare2012}. For the present paper, these histories matter because they help explain why a red button can feel eventful before any experiment begins.

Affordance theory sharpens that point but does not fully exhaust it. Gaver describes affordances as actionable properties that depend on both environment and perceiver \citep{Gaver1991}; Norman argues that design conventions stabilize which actions people believe are available \citep{Norman1999}. Experience-centered HCI adds that interaction is interpreted through felt meaning rather than through control logic alone \citep{McCarthyWright2004}. Pressability builds on these ideas by naming the relation that arises when symbol, invitation, bodily reach, and anticipated consequence align around a single object.

\subsection{Curiosity, boredom, agency, and observed action}

Voluntary pressing under minimal task demand sits close to research on curiosity, boredom, and agency. Curiosity is often triggered by uncertainty, novelty, and informational gaps \citep{KiddHayden2015}; boredom and under-stimulation can also drive people toward self-generated events, even when the activity is not obviously useful \citep{Wilson2014}. The button is therefore a useful HCI object precisely because it offers a clean, countable action without collapsing action into instrumental necessity.

This study also treats pressing as socially situated. Observable systems often change behavior simply by making action legible to self and others, a pattern that social facilitation research helps frame even when the setting is not explicitly interpersonal \citep{HallHenningsen2008}. The visible counter above the button extends that logic: every press becomes not just a motor response but a public contribution to an accumulating record. Agency here is not only the feeling that one can act, but the feeling that one's action enters a system and matters within it \citep{HaggardChambon2012}.

\subsection{Redness, salience, embodiment, and peripersonal space}

The button's perceptual design draws on work in color psychology and visual salience. Red frequently functions as an arousal and warning cue in human environments, but those meanings remain culturally contingent rather than universal \citep{ElliotMaier2014,GaoEtAl2007}. Salience research further shows that local feature contrast and singleton structure can attract attention even before reflective interpretation catches up \citep{Li2008,CastellottiEtAl2022}. In that sense, the big red button is not just culturally loud; it is engineered to become perceptually difficult to miss.

In immersive systems, salience is inseparable from bodily relation. Object placement in the visual field changes comfort, visibility, and selection efficiency \citep{Lin2021,Mirbagheri2024,Alghofaili2019,McNamara2019}. Work on embodiment and peripersonal space suggests that the felt boundary between self and environment is plastic, action-oriented, and sensitive to multisensory coupling \citep{KilteniEtAl2012,Serino2019,SerinoEtAl2018,BrozzoliEtAl2010}. This is especially relevant for the current study because the button is positioned not as scenery but as a reachable, body-adjacent object whose meaning may shift when it pulses with physiological signals.

\subsection{Biofeedback and physiological coupling in virtual reality}

Biofeedback in VR has been used for relaxation, self-regulation, and affective modulation, but the literature also shows that outcomes depend on how feedback is mapped, interpreted, and controlled \citep{LueddeckeFelnhofer2022,RockstrohEtAl2019}. Placebo control is particularly important because users can respond to the mere existence of a feedback loop, not only to its fidelity. Chittaro and colleagues demonstrate this directly in a breathing-focused VR comparison between real and placebo biofeedback \citep{ChittaroEtAl2024}. That design precedent matters here because the BRB study includes live and sham ECG conditions rather than assuming that any physiological aesthetic will automatically produce embodiment.

The current manuscript therefore treats physiological coupling as a tested proposition, not a built-in truth. If heartbeat-linked pulsing and respiration-linked lighting alter pressing, presence, or felt nearness to the button, that would support a stronger account of embodied pressability. If they do not, the null would still be informative: it would suggest that cultural history, salience, and ergonomic invitation may already do most of the explanatory work before physiology is added.
